This thesis is situated at the intersection of choreographic programming and automated state-space generation for probabilistic verification. In the following, we briefly discuss the most relevant lines of related work and position our contribution with respect to them.

\section{Choreographic Languages and Frameworks}
Montesi’s PhD thesis established key foundations of choreographic programming by introducing a formal model and type theory for the paradigm~\cite{Montesi2013}. As part of this research line, he also presented \textit{Chor}, a prototype framework that supports the development of choreography-based programs together with protocol validation and automatic Endpoint Projection (EPP) to executable Jolie code. By enabling developers to specify communication behavior from a single global perspective and verify it against protocol specifications, this approach supports the construction of distributed software that is safe by construction. Montesi’s later monograph systematizes these ideas and provides a broader overview of choreographic languages and their theoretical foundations~\cite{Montesi2023}. 

Building on these foundations, subsequent work proposed languages and frameworks such as \textit{Choral}~\cite{GiallorenzoMP24}, \textit{HasChor}~\cite{Shen_2023}, and \textit{AIOCJ}~\cite{lmcs:3263}, which explore object-oriented, functional, and adaptive variants of choreographic programming. \textit{Choral} reconstructs choreographic programming through object-oriented abstractions, modelling choreographies as classes and supporting distributed data types, modular composition, and compilation to pure Java libraries. \textit{HasChor}, in contrast, presents a functional formulation in which choreographies are represented as monadic computations in Haskell, with support for higher-order and location-polymorphic choreographies. \textit{AIOCJ} focuses on adaptable distributed applications by allowing designated parts of a choreography to be updated at runtime, while preserving key correctness guarantees such as deadlock freedom and race freedom.

In contrast to existing choreographic programming approaches, which predominantly focus on synthesizing executable distributed implementations from global specifications, this thesis introduced \textsc{ARIA} as a framework centered on formal model construction. Rather than targeting code generation, \textsc{ARIA} supports the automated derivation of analyzable stochastic models from high-level choreographies. In particular, it enables the construction of explicit-state Markov chain models with probabilistic choices and, where applicable, user-provided stochastic parameters. Such models can then be analyzed with probabilistic model checkers such as PRISM in order to study quantitative properties, including performance and reliability. In this way, \textsc{ARIA} extends choreography-based modeling towards quantitative formal analysis.

\section{Model Generation and State-Space Exploration}

Recent work has advanced choreographic programming by mechanizing its theory and compiler infrastructure. Cruz-Filipe et al.\ formalize the paradigm in \textit{Coq} and provide machine-checked proofs of key correctness properties, including deadlock freedom and the correctness of endpoint projection~\cite{cruzfilipe2021}. Building on this line of work, \textit{hacc}~\cite{cruzfilipe2023hacc} is the first certified compiler for a choreographic language, projecting choreographies from a core calculus to an intermediate process representation and ultimately to Jolie code. Similar efforts, such as \textit{Kalas}~\cite{pohjola_et_al2022}, likewise target verified end-to-end compilation from choreographies to executable distributed programs.

While these approaches focus on the translation of choreographic specifications into process code, \textsc{ARIA} considers a different direction, namely the construction of explicit-state stochastic models from choreography semantics. In doing so, it follows an on-the-fly explicit-state exploration paradigm comparable to that used in established model checkers such as SPIN~\cite{Holzmann03} and the explicit-state engine of PRISM~\cite{Prism}. Rather than relying on an intermediate process representation, \textsc{ARIA} explores the global state space by iteratively applying the operational semantics of the probabilistic choreography. This design supports the use of state deduplication and worklist-based traversal techniques during model construction.

\section{Probabilistic Verification of Distributed Systems}

The formal analysis of distributed systems through probabilistic model checking has a long tradition, with tools such as PRISM~\cite{Prism} and Storm~\cite{Storm} representing the current standard for the analysis of Discrete-Time and Continuous-Time Markov Chains (DTMCs and CTMCs). In conventional workflows, however, applying such tools typically requires the manual construction of a system model in a dedicated modeling language, such as the PRISM language. This process is often labor-intensive and error-prone, as it requires the modeller to encode the local behavior of each participant while ensuring that the resulting model correctly reflects the intended global synchronization.

To the best of our knowledge, comparatively less attention has been paid to the automated derivation of stochastic models for quantitative analysis from choreographic specifications. In this context, \textsc{ARIA} explores the construction of explicit-state Markov models from probabilistic choreographies, thereby providing a basis for quantitative formal analysis from a global specification.


